\chapter{System Testing}

\vspace*{3cm} 
\begin{center}
    {\fontsize{50}{50}\selectfont\bfseries Chapter Six:}
    
    \vspace{3.5cm}
    
    {\fontsize{30}{30}\selectfont\bfseries System Testing}
\end{center}

\thispagestyle{plain}



\vspace{1cm}

\begin{flushleft}

\noindent
5.1 Required Third-Party Tools
\dotfill \pageref{sec:tools}

\vspace{0.5cm}

\noindent
5.2 Installation Guide
\dotfill \pageref{sec:installation}

\vspace{0.5cm}

\noindent
5.3 Running the Application
\dotfill \pageref{sec:running}

\vspace{0.5cm}

\noindent
5.4 User Interface Walkthrough
\dotfill \pageref{sec:ui}

\end{flushleft}

\newpage


% ----------------------------------------------------------------
\section{Required Third-Party Tools}
\label{sec:tools}
% ----------------------------------------------------------------

\large

The following tools and libraries are required to run Revfit locally:

\begin{itemize}
    \item \textbf{Python 3.10+} with pip package manager
    \item \textbf{Flutter SDK} for the cross-platform frontend
    \item \textbf{CUDA-compatible GPU} (optional, for accelerated inference)
    \item \textbf{Spoonacular API Key} (for nutritional meal data) \cite{spoonacular}
\end{itemize}

\normalsize

\begin{table}[H]
\centering
\begin{tabular}{|l|l|}
\hline
\textbf{Tool / Library} & \textbf{Purpose} \\
\hline
uvicorn       & FastAPI ASGI development server \\
PyTorch       & Deep learning inference (VideoMAE, X3D-M, TCN) \\
MediaPipe     & Real-time 3D pose landmark extraction \\
OpenCV        & Video stream processing and HUD rendering \\
FAISS         & Vector database for collaborative filtering \\
FastAPI       & Backend web API \\
Flutter       & Cross-platform frontend user interface \\
Spoonacular   & External nutritional recipe API \\
\hline
\end{tabular}
\caption{Required third-party dependencies}
\label{tab:tools}
\end{table}

% ----------------------------------------------------------------
\section{Installation Guide}
\label{sec:installation}
% ----------------------------------------------------------------

\subsection*{Step 1: Clone Repository}

\begin{codebox}
\begin{verbatim}
git clone https://github.com/[your-repo]/Revfit.git
cd Revfit
\end{verbatim}
\end{codebox}

\subsection*{Step 2: Install Backend Dependencies}

\begin{codebox}
\begin{verbatim}
cd backend
pip install -r requirements.txt
\end{verbatim}
\end{codebox}

\large
This installs: FastAPI, uvicorn, torch, torchvision, mediapipe, opencv-python, faiss-cpu, and transformers (VideoMAE, X3D-M).
\normalsize

\subsection*{Step 3: Install Frontend Dependencies}

\begin{codebox}
\begin{verbatim}
flutter pub get
\end{verbatim}
\end{codebox}

\subsection*{Step 4: Configure Environment Variables}

Create a \texttt{.env} file in the project root:

\begin{codebox}
\begin{verbatim}
SPOONACULAR_API_KEY=your_api_key_here
\end{verbatim}
\end{codebox}

\subsection*{Optional: GPU Setup (NVIDIA)}

\begin{codebox}
\begin{verbatim}
pip install torch torchvision \
    --index-url https://download.pytorch.org/whl/cu121
\end{verbatim}
\end{codebox}

% ----------------------------------------------------------------
\section{Running the Application}
\label{sec:running}
% ----------------------------------------------------------------

\subsection*{Step 1: Start the Backend Server}

\begin{codebox}
\begin{verbatim}
cd backend
uvicorn main:app --reload
\end{verbatim}
\end{codebox}

\subsection*{Step 2: Start the Frontend}

\begin{codebox}
\begin{verbatim}
flutter run
\end{verbatim}
\end{codebox}

\subsection*{Step 3: Access the Application}

\large
The Flutter application will launch in the connected Android device or emulator.

\vspace{0.4cm}

Available API endpoints:
\begin{itemize}
    \item \url{http://127.0.0.1:8000/docs} --- Interactive Swagger UI documentation
    \item \url{http://127.0.0.1:8000/api/recommend} --- Recommendation engine endpoint
    \item \url{http://127.0.0.1:8000/api/analyze} --- Computer vision inference endpoint
\end{itemize}
\normalsize

\subsection*{Troubleshooting}

\begin{itemize}
    \item \textbf{Port busy}: run with \texttt{--port 8001} flag
    \item \textbf{CUDA error}: force CPU with \texttt{export CUDA\_VISIBLE\_DEVICES=-1}
    \item \textbf{Camera not detected}: ensure app permissions allow camera access on the Android device
\end{itemize}

\large
\noindent\textbf{Production deployment}: Remove \texttt{--reload}, build the frontend with \texttt{flutter build apk} to distribute to Android devices, and serve the backend via a reverse proxy (e.g., nginx).
\normalsize

% ----------------------------------------------------------------
\section{User Interface Walkthrough}
\label{sec:ui}
% ----------------------------------------------------------------

\subsection{Authentication and Onboarding}

\large

Upon launching the application, users are greeted with the Authentication screen. Following secure registration or login, the Onboarding flow collects the inputs required to construct the 30-dimensional user feature vector for the VDB. Users select their primary fitness goal, fitness level, available equipment, and dietary constraints.

\normalsize

% [INSERT FIGURE: Screenshot of Registration / Login Screen]


\vspace{0.3cm}

% [INSERT FIGURE: Screenshot of Onboarding Questionnaire]


\subsection{Personalized Dashboard}

\large

Once onboarding is complete, the Dashboard immediately populates with a personalized daily itinerary divided into the \textbf{Workout Routine} and \textbf{Nutritional Plan} segments, generated by the VDB-based CF engine.

\normalsize

% [INSERT FIGURE: Screenshot of User Dashboard]



\subsection{Real-Time Exercise Tracking}

\large

When a user initiates an exercise session, the Real-Time Tracking interface activates. The HUD displays:
\begin{itemize}
    \item A live repetition counter.
    \item A form status indicator (GOOD / VIOLATION).
    \item The current processing frame rate.
    \item A color-coded skeletal overlay (green = correct form; red = heuristic violation with directional text warning).
\end{itemize}

\normalsize

% [INSERT FIGURE: Screenshot of Real-Time Tracking — Correct Form]


\vspace{0.3cm}

% [INSERT FIGURE: Screenshot of Real-Time Tracking — Form Violation]


\subsection{Post-Workout Analysis Report}

\large

Upon set completion, the system aggregates timestamped kinematic data and presents a \textbf{Post-Workout Analysis Report} detailing:
\begin{itemize}
    \item Total valid repetitions counted.
    \item An itemized list of detected heuristic form violations, grouped by contiguous violation frames and presented as average angular deviations.
\end{itemize}

\normalsize

% [INSERT FIGURE: Screenshot of Post-Workout Analysis Report]